\section{序論}
11月13日朝の観測窓では、生成AIを実際の作業に組み込む製品発表が目立った。GoogleのGemini 4 Hydrogenは高速なコーディングや日常的なエージェント処理を志向し、AnthropicのClaude Meetingsは会議中の共同作業へAIの利用範囲を広げる。性能だけでなく、処理速度、価格、導入先の広さが比較の焦点となっている。

一方、NVIDIAの経営トップ交代を巡る報道がAI計算基盤の先行きへの関心を集めた。ゲノム解析モデルや判断専用APIへの言及も続き、AnthropicのIPO関連情報も観測された。以下ではGrokが収集した6件を記載順に整理する。独立した一次資料の再検証は行っていない。